\chapter{P5：先用案例决定要不要增加检索复杂度}

\section{承上启下}

P4 已经有了可重复查询的证据工具，但“更复杂”不等于“更好”。P5 要用同一批运营问题比较不同检索办法，回答：增加 BM25 或其他组件，是否带来值得维护的实际收益？

\section{业务问题、输入与边界}

输入是 P4 证据工具、P3 规则、同一商品范围，以及调参前固定的 27 个 silver 案例。业务价值是控制技术投入，不因为向量库或新名词看起来先进就增加部署、延迟和解释成本；技术价值是学会用固定案例和指标比较方案。

这些案例适合工程回归，不是独立人工 gold，也不代表所有开放问题。P5 不能据此声称真实用户已经更满意或生产质量已经达标。

\section{开发过程}

\subsection{1. 先把考题固定}

案例覆盖单品诊断、竞品比较、问题分流、商品歧义和证据不足。先固定考题，再运行词法基线和 FTS5/BM25，避免看到结果后修改问题。

\subsection{2. 用同一把尺子比较}

两个检索器消费同一数据版本、证据合同、商品范围和引用规则。除了命中，还看引用是否完整、商品是否关联正确、没有证据时是否会误报，以及延迟是否可接受。\textbf{Recall@k} 可以先理解为“前 k 个结果里找回了多少预期证据”。

\subsection{3. 接受没有胜出的结果}

当前 BM25 没有在冻结案例上胜出，所以没有为了制造技术亮点继续加入组件。还要注意：silver 期望本来由词法规则产生，用同一词法规则检索得到高命中，存在循环偏向。

\section{本阶段学到的知识}

\begin{itemize}
  \item 冻结案例：先把考题和答案边界固定。
  \item case hit：一个案例是否找到了它预期的证据。
  \item Recall@k：前 k 个结果找回预期证据的比例。
  \item 引用完整性：答案里的每个引用能否回到实际证据。
  \item 无答案误报：没有证据时，系统是否硬写结论。
  \item 延迟分位数：多数请求和较慢请求分别需要多久。
\end{itemize}

\begin{beginnerbox}
如果“标准答案”本来就是用某个词表产生的，那么用同一个词表检索到高命中，只能说明规则前后一致，不能证明系统真正理解了语义。
\end{beginnerbox}

\section{问题、修复与验证}

问题一是把技术复杂度当成先进程度；用冻结案例和统一指标修复。问题二是结果不理想时反复改案例；修复是保留原始案例和 Bad Case。问题三是把 silver 回归分数当成人工质量；报告中明确写出标签循环偏向。验证结果只支持一个谨慎结论：在当前范围和案例上，没有证据说明增加复杂度值得。

\section{阶段产出与验收边界}

公开产出包括\filelink{artifacts/evaluation-summaries/phase-5-retrieval.json}和\filelink{docs/reports/phase-5-retrieval-evaluation.md}。P5 在其范围内为\statuspassed：比较、引用、关联、无答案和延迟门通过；不表示词法检索已经达到生产质量。

\section{怎样进入下一阶段}

P6 要解决 P5 留下的业务缺口：运营人员还不能用自然语言完成“选商品、取证、组织结果、澄清歧义和安全失败”这一整条任务。只有当这个缺口出现后，加入 Agent 才有明确理由。

\section{面试怎么讲}

\begin{interviewbox}
“我用冻结案例比较检索方案，而不是凭感觉引入向量库。BM25 没有胜出，我保留了负结果；同时说明 silver 标签由词法规则产生，存在循环偏向。因此 P5 只支持‘当前没有证据需要增加复杂度’，不把回归分数说成真实语义质量。”
\end{interviewbox}

\textbf{掌握标准：}能解释为什么先固定案例、为什么保留负结果，以及为什么 P5 通过仍不足以证明业务收益。
